\section{Thermal Treatment and Densification}

\subsection{Pressureless schedules}

Pressureless firing is not a single baseline. The audited library contains a title-level study of pressure filtration followed by sintering for Al-containing LLZTO \cite{lei2021fabrication}; it marks shaping plus firing but provides no locally checked density or optical-property value. A stored abstract systematically varies lithium excess and dwell time in Ga-containing LLZO and reports that these parameters matter within its processing case \cite{schwab2023towards}. Another abstract studies oxygen-assisted firing and reports void, density, conductivity, hardness, and cell-related readouts \cite{shen2020a}. Taken together, these records show that schedule, lithium inventory, atmosphere, and green-body preparation must remain visible even when pressureless firing is the primary tag.

Uniformity adds a scale within the specimen. The stored abstract of a powder-covering strategy considers particle-size distribution, heating rate, support geometry, densification evolution, and spatially uniform readouts \cite{guo2023uniform}. This is evidence that uniformity can be a measured target rather than an assumed consequence of average density. It does not establish that the reported strategy transfers unchanged to every composition or furnace.

\subsection{Pressure, field, and ultrafast assistance}

The title-level corpus contains papers explicitly framed around spark-plasma-sintering mechanisms \cite{yamada2016sintering}, field-assisted densification of Al-substituted cubic LLZO \cite{zhang2014field}, rapid ultra-high-temperature processing of Ga-containing LLZO \cite{tao2022preparation}, and ultrafast high-temperature sintering of LLZO \cite{ihrig2021li7la3zr2o12}. These keys establish that pressure, field, and ultrafast branches are represented. They do not support energy-efficiency, mechanism, density, or conductivity rankings in this manuscript.

Two abstracts permit bounded comparisons. One compares solid-state sintering, SPS, and SPS followed by heat treatment using phase, morphology, and ionic-transport readouts \cite{xue2020spark}. Another varies hot-pressing temperature at a reported fixed pressure and reports relative density, grain-boundary resistance, and room-temperature total conductivity \cite{david2015microstructure}. Each is informative within its own specimens. Across studies, however, starting powder, dopant content, pressure path, ramp, peak temperature, dwell, post-treatment, and measurement temperature differ or may be unavailable. We therefore do not rank ``SPS,'' ``hot press,'' ``field assisted,'' and ``ultrafast'' as universal technologies.

\subsection{Low-temperature and additive-assisted densification}

``Low temperature'' is meaningful only relative to a comparator and material system stated in the same source. A cold-sintering abstract describes an LLZO-based composite, a transient liquid, microstructure, conductivity, and a cell demonstration \cite{seo2020broad}. Because this is a composite route, it is not silently normalized to monolithic LLZO firing. A SiO2-addition abstract reports removal of surface carbonate, densification, grain-boundary transport, and cell readouts, while also advancing a mechanistic interpretation \cite{zhang2020enhanced}. A CuO study reports a lower firing schedule for W-containing garnet together with microstructural, transport, wetting, and cell measurements \cite{zhang2019effects}. A 2024 Li5AlO4-addition study reports a within-paper comparison involving firing temperature, relative density, conductivity, critical current, and cell operation \cite{zhou2024li5alo4}. These are associated outcomes from abstract-bearing records, not proof of one common densification pathway.

The evidence rule is therefore to record achieved density separately from the mechanism proposed to explain it. An additive may supply lithium, alter a transient phase, change boundary chemistry, or correlate with particle rearrangement; an abstract-level association cannot by itself discriminate those alternatives. A fair comparison must preserve parent powder, additive dose, green density, pressure, complete thermal history, final phase assemblage, density basis, and separated transport if available. Otherwise, ``lower temperature'' remains a study-specific descriptor rather than a league table.
